\documentclass[10pt,a4paper]{article}

% -------------------------------------------------
% Page layout
% -------------------------------------------------
\usepackage[
    left=2.15cm,
    right=2.15cm,
    top=2.10cm,
    bottom=1.90cm
]{geometry}

\usepackage[T1]{fontenc}
\usepackage{lmodern}
\usepackage{xcolor}
\usepackage{listings}
\usepackage{fancyhdr}
\usepackage{titlesec}
\usepackage{microtype}

% -------------------------------------------------
% Colours approximating the uploaded document
% -------------------------------------------------
\definecolor{codebackground}{RGB}{238,237,252}
\definecolor{codeblue}{RGB}{35,42,145}
\definecolor{rulegray}{RGB}{80,80,80}

% -------------------------------------------------
% Running header and footer
% -------------------------------------------------
\pagestyle{fancy}
\fancyhf{}

\fancyhead[L]{\footnotesize 3,230: Data Handling: Import, Cleaning and Visualisation}
\fancyhead[R]{\footnotesize Autumn 2026: Summary by Leroy Zahno}
\fancyfoot[C]{\footnotesize \thepage}

\renewcommand{\headrulewidth}{0.35pt}
\renewcommand{\footrulewidth}{0.35pt}

\setlength{\headheight}{33pt}
\setlength{\headsep}{12pt}
\setlength{\footskip}{8pt}

% -------------------------------------------------
% Heading styles
% -------------------------------------------------
\titleformat{\section}
  {\normalfont\large\bfseries}
  {}
  {0pt}
  {}

\titlespacing*{\section}
  {0pt}
  {0pt}
  {0.8em}

\titleformat{\subsubsection}
  {\normalfont\normalsize\bfseries}
  {}
  {0pt}
  {}

\titlespacing*{\subsubsection}
  {0pt}
  {0pt}
  {0.8em}


\titleformat{\subsection}
  {\normalfont\normalsize\bfseries}
  {}
  {0pt}
  {}

\titlespacing*{\subsection}
  {0pt}
  {1.1em}
  {0.3em}

% -------------------------------------------------
% R console style
% -------------------------------------------------
\lstdefinestyle{rconsole}{
    language=R,
    basicstyle=\ttfamily\normalsize,
    backgroundcolor=\color{codebackground},
    keywordstyle=\color{codeblue},
    commentstyle=\color{codeblue},
    stringstyle=\color{codeblue},
    identifierstyle=\color{codeblue},
    numberstyle=\color{codeblue},
    columns=fullflexible,
    keepspaces=true,
    showstringspaces=false,
    breaklines=true,
    breakatwhitespace=false,
    frame=none,
    xleftmargin=0pt,
    xrightmargin=0pt,
    framexleftmargin=0pt,
    framexrightmargin=0pt,
    aboveskip=0.15em,
    belowskip=0.45em,
    lineskip=-0.5pt
}

% Make all console text blue, as in the source document
\lstset{
    style=rconsole,
    literate=
      {>}{{{\color{codeblue}>}}}1
      {[}{{{\color{codeblue}[}}}1
      {]}{{{\color{codeblue}]}}}1
}

% Slightly compact paragraph formatting
\setlength{\parindent}{0pt}
\setlength{\parskip}{0pt}

\begin{document}

% -----------------------------------
% START
% -----------------------------------

\section{General Functions}

\subsection*{\texttt{data()}}

This function loads a data set, which is often a built-in part of R.

 \vspace{.5em}

\begin{lstlisting}
> data(iris)
>
\end{lstlisting}
 This command loads the data.frame "iris" into your RAM and then lets you type your next command into the terminal. There will be no further feedback from the shell as is often usual with other commands.

\subsection*{\texttt{library()}}

This function makes an add-on package available.

 \vspace{.5em}

\begin{lstlisting}
> library(stringr)
> library(dplyr)
> library(tidyverse) # This library contains ggplot2
>
\end{lstlisting}
 This command loads the data.frame "iris" into your RAM and then lets you type your next command into the terminal. There will be no further feedback from the shell as is often usual with other commands.
 
 \vspace{1em}

\section{Data Import}

\subsection*{\texttt{read\_csv()}}

\begin{lstlisting}
# You can either browse to the desired folder as shown below, or setwd()
# to the relevant directory and directly type in the file name, without apostrophes.

> read_csv("data/file1.csv")
> 
\end{lstlisting}

\section{Data Cleaning}

\subsection*{\texttt{str\_split()}}

\begin{lstlisting}
> library(stringr)
> str_split("banana", pattern= "a")
[[1]]
[1] "b" "n" "n" "" 
\end{lstlisting}

\subsection*{\texttt{filter(data.frame(\#column\_name = c(), \#repeat\_n\_times), range)}}

\begin{lstlisting}
# The last row rounds values in the numbers column to the second decimal digit and replaces the decimal mark 

> library(dplyr)
>

RScript:
df <- data.frame(id = c(1:3,1:3),
                 numbers= c(1.456, 5.784356, 3.123, 8.234, 23.5465, 2.345),
                 year=c(2017,2017,2017,2018,2018,2018))
df_filtered <- filter(df, numbers < 5 & year== 2017)
numbers_formatted <- format(round(df_filtered$numbers, 2), decimal.mark=",")
\end{lstlisting}

\subsection*{\texttt{summarise(filter(data.frame(\#column\_name = c(), \#repeat\_n\_times), range)}}

\begin{lstlisting}
df_c <- data.frame(id = c(1:3,1:3),
                   money_spent= c(1000, 2000, 6000, 1500, 3000, 5500),
                   currency = c("CHF", "CHF", "USD", "EUR", "CHF", "USD"),
                   year=c(2017,2017,2017,2018,2018,2018))
df_filtered <- filter(df_c, currency== "CHF" )
df_summary <- summarise(df_filtered,
                        total = sum(money_spent),
                        N_obs = n())
\end{lstlisting}
			
\vspace{1em}

\subsection{mutate()}

The mutate() function in R is used to create new variables or modify existing variables in a data frame without removing any other variables. It allows you to apply transformations or calculations to columns and add the result as new columns or overwrite existing ones.

\vspace{0.5em}

\begin{lstlisting}
library(dplyr) 

d <- data.frame( name = c("Abhi", "Bhavesh", "Chaman", "Dimri"), 
                 age = c(7, 5, 9, 16), 
                 ht = c(46, NA, NA, 69), 
                 school = c("yes", "yes", "no", "no") ) 

mutate_d <- mutate(d, x3 = ht + age) 

\end{lstlisting}

\vspace{1em}

\section{Data Visualisation}

\subsection*{\texttt{ggplot()}}

Every ggplot2 plot has at least three key components: data, aes, and a geom function (or another rendering-instruction):

\vspace{0.5em}

\begin{lstlisting}
> data(iris)
> library(ggplot2)
> ggplot(iris, aes(x=Sepal.Length, y=Sepal.Width)) + geom_point()
>
\end{lstlisting}

\vspace{1em}

\subsubsection{Typical Errors}

$\textbf{Error 1 (Lack of attributes):} \quad$ You can see that the third attribute is missing, which is why there will be a plot, however it is going to be empty:

\vspace{0.5em}

\begin{lstlisting}
> data("swiss")
> library(tidyverse)
> ggplot(swiss, aes(x=Agriculture, y=Examination))
\end{lstlisting}

\vspace{0.5em}

$\textbf{Error 2 (Faulty configuration of functions):} \quad$ This is also not going to work for several reasons: (1) The data.frame has not been loaded into R, (2) since there is no df-object, the axes "A" and "B" will also have no meaning to the aes-attribute within the ggplot-function, and (3) there is no color "C" in the geom\_point function's aesthetics.

\vspace{0.5em}

\begin{lstlisting}
> ggplot(df, aes(x=A, y=B)) + geom_point(aes(colour=C)) + geom_smooth(colour="red", method = "lm")
>
\end{lstlisting}

\vspace{0.5em}

Assume that there is an object called "df", then:

\begin{enumerate}
\item The line geom\_smooth(colour="red", method = "lm") adds an additional layer to the plot, displayed
on top of what geom\_point(aes(colour=C)) generates.
\item df must be a data.frame, tibble, or data.table with columns A, B, and C in order to make this code
work (result in a plot).
\item The resulting plot will have a legend.
\end{enumerate}

\end{document}